\chapter{Radio Wave Propagation}
\label{ch:propagation}

How a radio signal gets from a transmitting antenna to a receiving
antenna, hundreds or thousands of kilometres away, depends on frequency,
time of day, season, solar activity, and geography. This chapter surveys
the main propagation mechanisms an amateur exploits across the HF, VHF,
UHF, and microwave amateur bands.

\section{The Ground Wave}

At lower frequencies (roughly below a few MHz), a vertically polarized
wave can travel along the Earth's curved surface for a considerable
distance as a \textbf{ground wave} (or \textbf{surface wave}), guided
partly by the interaction between the wave and the conductive ground.
Ground-wave range is limited by ground conductivity (seawater performs
far better than dry rocky or sandy soil) and increases at lower
frequency; it is a minor factor for most HF amateur bands but is the
dominant propagation mode for AM broadcast and some maritime/beacon
services at lower frequencies.

\section{Line-of-Sight and the Radio Horizon}

At VHF, UHF, and microwave frequencies, propagation is normally
\textbf{line-of-sight}: the direct wave travels essentially in a
straight line, blocked by the curvature of the Earth and by terrain
obstructions. Because the atmosphere's refractive index normally
decreases slightly with height, radio waves bend (\textbf{refract})
gently downward, extending the effective \textbf{radio horizon} somewhat
beyond the optical horizon --- a commonly used approximation for the
radio horizon distance (in kilometres, with antenna height $h$ in
metres) is
\[
d_{\text{km}} \approx 4.12\sqrt{h_{\text{m}}}
\]
applied to each end of the path and summed, which is why raising a
VHF/UHF antenna (or operating from a hilltop) so directly and
dramatically increases usable range.

\subsection{Tropospheric Effects}

Under certain weather conditions, temperature and humidity gradients in
the lower atmosphere (\textbf{troposphere}) create layers where the
refractive index changes more sharply than usual, bending
(\textbf{tropospheric ducting} or simple \textbf{tropo enhancement}) VHF
and UHF signals well beyond the normal radio horizon --- sometimes
hundreds of kilometres --- along stable high-pressure weather boundaries
or temperature inversions. \textbf{Tropospheric scatter}, a weaker but
more consistently available mechanism, allows some signal energy to
scatter forward off small-scale turbulence in the troposphere even
without a defined duct.

\section{The Ionosphere and Skywave Propagation}

The \textbf{ionosphere} is a region of the upper atmosphere
(approximately 60--1000\,km altitude) where solar ultraviolet and X-ray
radiation ionizes atmospheric gases, creating layers of free electrons
capable of refracting (and, at a steep enough angle or low enough
frequency, fully reflecting) radio waves back toward Earth. This
\textbf{skywave} propagation is what makes long-distance (``DX'') HF
communication possible without a chain of relay stations, by bouncing a
signal off the ionosphere and back down to a distant receiver --
potentially over several successive ``hops.''

\subsection{Ionospheric Layers}

\begin{itemize}
\item \textbf{D layer} (${\sim}60$--90\,km): forms only in daylight;
does not refract HF signals back to Earth but does \emph{absorb} them,
particularly at lower HF frequencies --- the main reason low HF bands
(80\,m, 40\,m) are noticeably ``quieter'' but shorter-range for DX
during the day and open up dramatically for longer-distance
propagation after D-layer ionization decays at night.
\item \textbf{E layer} (${\sim}90$--150\,km): present mainly in
daylight, refracts lower/mid HF frequencies and supports the anomalous,
short-lived \textbf{sporadic-E} propagation discussed below.
\item \textbf{F layer} (splitting into \textbf{F1} and \textbf{F2} in
daylight, recombining into a single F layer at night, roughly
150--500\,km): the primary layer responsible for long-distance HF
skywave propagation, with the F2 layer's height and ionization density
--- strongly dependent on solar activity --- setting the
\textbf{maximum usable frequency (MUF)} for a given path on a given day.
\end{itemize}

\subsection{Critical Frequency and MUF}

For a signal transmitted straight up, the highest frequency the
ionosphere will still reflect back down is the \textbf{critical
frequency}, $f_c$. For a signal transmitted at an oblique angle
(the normal case for a real communication path), a higher frequency
can still be refracted back to Earth --- the \textbf{maximum usable
frequency} for that specific path, related to the critical frequency and
the wave's incidence angle $\theta$ (measured from vertical) by the
\textbf{secant law}:
\[
\boxed{\text{MUF} = f_c \sec(\theta)}
\]
Operating above the MUF for a given path, the signal passes through the
ionosphere into space rather than being returned to Earth. Because lower
launch angles (more oblique, larger $\theta$) permit a higher MUF, DX
contacts over very long paths are often possible at frequencies that
would not support shorter, steeper-angle paths at the same time.

\begin{keyidea}[LUF, MUF, and the usable window]
Below a \textbf{lowest usable frequency (LUF)}, D-layer absorption (and
increasing atmospheric/man-made noise at lower frequencies) makes a
signal too weak to be useful even if the ionosphere would otherwise
support the path; above the MUF, the signal is not returned to Earth at
all. Real-world HF operation lives in the changing window between these
two limits, which is why band selection ``for the conditions'' is such
a core operating skill --- and why real-time MUF/LUF prediction tools
and propagation beacons are widely used by serious HF operators.
\end{keyidea}

\section{The Solar Cycle and Space Weather}

Ionization of the F2 layer depends strongly on solar ultraviolet
output, which varies on an approximately \textbf{11-year solar cycle} of
rising and falling sunspot activity. Higher sunspot numbers generally
mean a denser, higher F2 layer and a correspondingly higher MUF,
supporting DX propagation on higher HF bands (15\,m, 12\,m, 10\,m) that
may be largely closed during solar minimum. Solar flares and coronal
mass ejections can also cause sudden, disruptive space-weather events:

\begin{itemize}
\item A \textbf{sudden ionospheric disturbance (SID)}, triggered by a
solar flare's X-ray flux, causes a rapid increase in D-layer ionization
and a corresponding \textbf{HF radio blackout} on the sunlit side of
Earth, lasting from minutes to a few hours.
\item A \textbf{geomagnetic storm}, triggered by a coronal mass
ejection reaching Earth's magnetosphere, disturbs the ionosphere over a
longer period (a day or more), variably degrading HF propagation while
sometimes \emph{enhancing} high-latitude VHF/UHF propagation via aurora.
\end{itemize}

\section{Sporadic-E}

\textbf{Sporadic-E} propagation results from dense, patchy, short-lived
clouds of ionization in the E layer (mechanism not fully understood, but
statistically linked to season --- most common in late spring/early
summer in each hemisphere --- and to wind shear in the upper
atmosphere). It can support surprisingly strong single- or
multi-hop propagation on the upper HF bands and, notably, on the lower
VHF bands (6\,m and, less commonly, 2\,m) that are normally limited to
line-of-sight --- producing some of the most exciting unpredictable
VHF DX openings available to an amateur.

\section{Meteor Scatter and Auroral Propagation}

\begin{itemize}
\item \textbf{Meteor scatter}: the brief ionized trail left by a meteor
burning up in the upper atmosphere can reflect VHF signals for a
fraction of a second to a few seconds, long enough (especially with
digital modes designed for very short contacts) to complete brief
contacts over 1000--2000\,km paths, enhanced during annual meteor
showers.
\item \textbf{Auroral propagation}: during geomagnetic disturbances,
VHF/UHF signals can scatter off ionization associated with the aurora,
typically producing a characteristic rasping, Doppler-smeared signal
quality, mainly usable for paths oriented toward the auroral region at
higher latitudes.
\end{itemize}

\section{Greyline Propagation}

The \textbf{greyline} (or \textbf{terminator}) is the band of twilight
circling the Earth between day and night. Along the greyline, the D
layer on the sunset/sunrise side has largely decayed (reducing
absorption) while the F layer, which decays more slowly after sunset,
remains sufficiently ionized to support propagation --- a favorable
combination that experienced HF operators specifically exploit for
long-path and unusual DX openings around their local sunrise and
sunset.

\section{Multipath, Fading, and Doppler Effects}

A received signal often arrives via more than one path simultaneously
(different ionospheric hops, ground wave plus skywave, or multiple
ionospheric layers at once), each with a slightly different path length
and therefore a different arrival phase. The combination of these
paths, constantly shifting as the ionosphere changes, produces
\textbf{fading} --- and, over sufficiently different path lengths,
measurable time-domain echo or \textbf{Doppler}-related frequency
smearing, both of particular relevance to digital mode performance
(Chapter~\ref{ch:digital-modes}) and to precise frequency measurement.

\begin{quickreview}
\item Ground wave (low HF/MF, vertical polarization) and line-of-sight
plus tropospheric effects (VHF/UHF/microwave) dominate at those
frequency extremes; skywave dominates general-purpose HF DX.
\item Ionospheric D, E, and F layers absorb (D) or refract (E, F)
HF signals; the F2 layer sets the MUF for long-distance skywave paths,
MUF $= f_c\sec\theta$.
\item The solar cycle and space-weather events (SIDs, geomagnetic
storms) drive large day-to-day and year-to-year variation in HF
propagation.
\item Sporadic-E, meteor scatter, and auroral propagation are
anomalous mechanisms that can open otherwise line-of-sight-limited
VHF bands over much longer paths.
\item Greyline propagation exploits the favorable D-layer/F-layer
balance found along the day/night terminator.
\end{quickreview}

\begin{practiceproblems}
\item Estimate the radio horizon distance for a VHF station with a
30\,m antenna height, communicating with a station whose antenna is
10\,m high (sum both contributions).
\item A path has a critical frequency of 4\,MHz and an incidence angle
of $70\degree$ from vertical. Estimate the MUF for this path.
\item Explain why the 80\,m and 40\,m amateur bands are typically much
better for long-distance contacts at night than during the day.
\item Explain, qualitatively, why higher sunspot activity tends to
improve propagation on the higher HF bands (e.g., 10\,m) specifically.
\item Explain why meteor scatter contacts are typically brief and rely
on specialized digital modes rather than sustained voice conversation.
\end{practiceproblems}
